\documentclass[11pt,a4paper]{article}

% --- packages -------------------------------------------------------------
\usepackage[utf8]{inputenc}
\usepackage[T1]{fontenc}
\usepackage{lmodern}
\usepackage[margin=2.4cm]{geometry}
\usepackage{booktabs}
\usepackage{array}
\usepackage{xcolor}
\usepackage{enumitem}
\usepackage{titlesec}
\usepackage{fancyhdr}
\usepackage{listings}
\usepackage[hidelinks]{hyperref}

% --- colours & styling ----------------------------------------------------
\definecolor{accent}{HTML}{1F4E79}
\definecolor{ok}{HTML}{1B7F3B}
\definecolor{bad}{HTML}{B00020}
\definecolor{codebg}{HTML}{F4F6F8}
\definecolor{rule}{HTML}{C9D3DD}

\titleformat{\section}{\Large\bfseries\color{accent}}{\thesection}{0.6em}{}
\titleformat{\subsection}{\large\bfseries\color{accent}}{\thesubsection}{0.6em}{}
\setlist[itemize]{leftmargin=1.3em,itemsep=2pt,topsep=2pt}

\pagestyle{fancy}
\fancyhf{}
\lhead{\small\color{accent}Linux Hardening Auditor}
\rhead{\small\color{accent}Project Closeout Report}
\cfoot{\small\thepage}
\renewcommand{\headrulewidth}{0.4pt}

\lstdefinestyle{console}{
  basicstyle=\ttfamily\footnotesize,
  backgroundcolor=\color{codebg},
  frame=single,
  rulecolor=\color{rule},
  framesep=6pt,
  breaklines=true,
  columns=fullflexible,
  keepspaces=true,
  showstringspaces=false,
}

\newcommand{\pass}{\textcolor{ok}{\textbf{PASS}}}
\newcommand{\fail}{\textcolor{bad}{\textbf{FAIL}}}

% --- title ----------------------------------------------------------------
\title{\vspace{-1.2cm}\color{accent}\Huge\textbf{Linux Hardening Auditor}\\[4pt]
\Large A CIS-mapped audit \& safe-remediation tool for Ubuntu 22.04\\[10pt]
\large\normalcolor Project Closeout Report}
\author{Milad Bahari Qaragoz}
\date{24 June 2026}

\begin{document}
\maketitle
\thispagestyle{fancy}

% =========================================================================
\section{Executive summary}
The Linux Hardening Auditor is a Python tool that audits a Linux host against the
\textbf{CIS Ubuntu 22.04 LTS Benchmark v1.0.0} (with BSI IT-Grundschutz alignment where
relevant), reports pass/fail per control with severity and remediation, and can optionally and
\emph{reversibly} fix what it finds. It is read-only by default; every mutating action backs up
the file it touches first and re-verifies the control afterwards.

The tool was validated end-to-end on a stock \textbf{Ubuntu 22.04.5 LTS} cloud host. A default
image scored \textbf{53\,\%}; after a single \texttt{fix} run it scored \textbf{100\,\%}, with
all 15 remediable controls remediated \emph{and} independently re-verified against the live host.

\begin{center}
\renewcommand{\arraystretch}{1.3}
\begin{tabular}{l c c}
\toprule
\textbf{Stage} & \textbf{Score} & \textbf{Controls passing} \\
\midrule
Before --- default Ubuntu 22.04 image & \textcolor{bad}{\textbf{53\,\%}} & 17 / 32 \\
After --- single \texttt{fix} run, re-audited & \textcolor{ok}{\textbf{100\,\%}} & 32 / 32 \\
\bottomrule
\end{tabular}
\end{center}

\paragraph{How to read the score.} These percentages are measured \emph{against this tool's own
catalogue of 32 checks}, not against the complete benchmark. ``100\,\%'' means ``all 32
implemented controls pass'' --- it is \textbf{not} a claim of full CIS compliance or of absolute
host security. The published CIS Ubuntu 22.04 Benchmark defines several hundred recommendations;
this project implements a deliberate, documented subset (\S4, and \texttt{docs/CONTROLS.md}). The
result should always be read alongside the catalogue that defines it: it states that the controls
\emph{this tool knows how to check} are satisfied, and nothing beyond that.

% =========================================================================
\section{Design principles}
The codebase is judged against three non-negotiable principles, chosen so the project reads like
the work of someone who builds auditable, maintainable security tooling:

\begin{itemize}
  \item \textbf{Scalable.} Controls are data-driven, self-registering plugins. Adding control
        \#33 means writing one self-contained check module and nothing else --- the engine,
        reporters, and remediator never grow per control.
  \item \textbf{Traceable.} Every control maps to a real CIS control ID. A reader can go from a
        CV bullet $\rightarrow$ report line $\rightarrow$ control ID $\rightarrow$ the exact
        check code $\rightarrow$ the written rationale, with no missing link. Six Architecture
        Decision Records capture \emph{why}, not just what.
  \item \textbf{Debuggable.} Structured logging at every check, deterministic output ordered by
        control ID, dry-run by default, and clear ``expected vs.\ found'' failure messages.
\end{itemize}

% =========================================================================
\section{Architecture}
The engine discovers and runs whatever checks are registered and emits a uniform result that the
reporters render; it never names a control.

\begin{lstlisting}[style=console]
auditor/
  cli.py          # entrypoint: audit | fix
  engine.py       # discovers + runs checks, collects results
  registry.py     # control + fixer discovery (@control / @fixer)
  models.py       # Control / CheckResult / Severity / Status
  host.py         # injected read-only host (LocalHost / FakeHost)
  reporters/      # console - markdown - html - json
  remediation.py  # backup -> dry-run -> apply -> verify
  checks/         # ONE self-registering module per control
\end{lstlisting}

\noindent\textbf{Adding a control} is a five-step, single-file path: create
\texttt{checks/<id>\_<slug>.py}, let it self-register via the \texttt{@control} decorator, add a
catalogue row and a threat-model rationale, add a test, add a changelog line. The host is
injected, so every check is unit-testable on a non-Linux dev box through a \texttt{FakeHost}.

% =========================================================================
\section{Control catalogue}
The catalogue covers \textbf{32 controls} spanning SSH hardening, kernel \texttt{sysctl} network
hardening, file permissions, password policy, the host firewall, auditing, and automatic
security updates.

\begin{center}
\renewcommand{\arraystretch}{1.25}
\begin{tabular}{l c p{8.2cm}}
\toprule
\textbf{Severity} & \textbf{Count} & \textbf{Examples} \\
\midrule
High   & 7  & SSH \texttt{PermitRootLogin no}, host firewall active, auditd enabled, \texttt{/etc/shadow} mode \\
Medium & 23 & \texttt{sysctl} network hardening, SSH MACs/KEX/MaxAuthTries, password min-length \& expiry \\
Low    & 2  & cron daemon active, password-expiry warning age \\
\midrule
\multicolumn{2}{l}{\textbf{Total: 32}} & \\
\bottomrule
\end{tabular}
\end{center}

% =========================================================================
\section{Remediation engine}
Remediation is split into three independently testable parts (ADR-0005):

\begin{enumerate}[leftmargin=1.5em,itemsep=2pt]
  \item A control's \textbf{fixer} is a pure function \texttt{fix(host) -> [Action]}; it only
        \emph{reads} the host and returns a declarative plan, so \texttt{--dry-run} is ``plan,
        then print'' with zero side effects.
  \item An \textbf{Applier} performs the side effects, \textbf{backing up every file before it
        changes it}. SSH changes are written as drop-in files and validated with
        \texttt{sshd -t} \emph{before} the daemon is reloaded, so a bad config can never lock
        you out.
  \item After a control's actions are applied, the control's \textbf{own check is re-run to
        verify} --- a fix that does not make the check pass is reported as failed, not assumed.
\end{enumerate}

\subsection{The shared-file bug and its fix (ADR-0006)}
Two controls legitimately edit the same file: \texttt{5.5.1.1} (\texttt{PASS\_MAX\_DAYS}) and
\texttt{5.5.1.2} (\texttt{PASS\_MIN\_DAYS}) both rewrite \texttt{/etc/login.defs}. Because each
fixer captured the file's content \emph{at plan time}, applying them in sequence let the second
whole-file write silently clobber the first fixer's change --- only the last writer of a shared
file survived. This surfaced during the cloud demo run.

\textbf{Fix:} \texttt{apply} now re-derives each control's actions by calling its fixer again
against the \emph{live} host at apply time, so every fixer sees the previous one's edits. The
backup layer additionally keeps the \emph{first} backup of a file touched twice in one run, so
the saved copy is the true pre-run original. The up-front plan still drives the read-only
\texttt{--dry-run} preview. Both the shared-file case and the command-failure abort path are
covered by regression tests.

% =========================================================================
\section{Real-host verification}
\subsection{Method}
The tool was run end-to-end on a freshly created \textbf{Ubuntu 22.04.5 LTS} Google Compute
Engine VM (\texttt{e2-micro}, Python 3.10.12) --- the exact distribution and interpreter the
benchmark targets. The audit read live \texttt{sshd -T} output, \texttt{/proc/sys},
\texttt{/etc/login.defs}, and real file modes; it was reached over an IAP tunnel with no public
IP, and given outbound-only internet via Cloud NAT so package-installing fixers could run.

\subsection{Before --- default image (53\,\%, 17/32)}
\begin{lstlisting}[style=console]
Linux Hardening Audit -- lha-audit-vm -- 2026-06-24T08:56:21
Score: 53%  (17 / 32 controls passing)

HIGH
  [FAIL] 3.5.1.3  Host firewall active        found "ufw inactive"
  [FAIL] 4.1.1.2  auditd enabled and active   found "active=inactive"
  [FAIL] 5.2.7    SSH PermitRootLogin disabled found "without-password"
MEDIUM
  [FAIL] 1.5.4    Core dumps restricted        fs.suid_dumpable=2 (want 0)
  [FAIL] 3.3      Kernel network hardening     1 of 7 not hardened
  [FAIL] 5.2.12   SSH X11Forwarding disabled   found "yes"
  [FAIL] 5.2.14   SSH strong MACs only         weak MACs present
  [FAIL] 5.2.16   SSH AllowTcpForwarding off   found "yes"
  [FAIL] 5.2.18   SSH MaxAuthTries <= 4        found "6"
  [FAIL] 5.2.21   SSH LoginGraceTime <= 60s    found "120"
  [FAIL] 5.4.1    Password min length >= 14    pwquality.conf missing
  [FAIL] 5.5.1.1  Password expiration <= 365   PASS_MAX_DAYS=99999
  [FAIL] 5.5.1.2  Min days between changes >=1 PASS_MIN_DAYS=0
  ... (15 failing in total)
\end{lstlisting}

\subsection{Apply --- 15 controls remediated and verified}
\begin{lstlisting}[style=console]
Remediation applied (backups in backups/20260624-085815/):
  [OK] 1.5.4   Core dumps restricted          -- now passing
  [OK] 3.3     Kernel network hardening        -- now passing
  [OK] 3.5.1.3 Host firewall active            -- now passing
  [OK] 4.1.1.2 auditd enabled and active       -- now passing
  [OK] 5.1.2   /etc/crontab permissions        -- now passing
  [OK] 5.2.1   sshd_config permissions         -- now passing
  [OK] 5.2.7   SSH PermitRootLogin disabled    -- now passing
  [OK] 5.2.12  SSH X11Forwarding disabled      -- now passing
  [OK] 5.2.14  SSH strong MACs only            -- now passing
  [OK] 5.2.16  SSH AllowTcpForwarding disabled -- now passing
  [OK] 5.2.18  SSH MaxAuthTries <= 4           -- now passing
  [OK] 5.2.21  SSH LoginGraceTime <= 60s       -- now passing
  [OK] 5.4.1   Password minimum length >= 14   -- now passing
  [OK] 5.5.1.1 Password expiration <= 365 days -- now passing
  [OK] 5.5.1.2 Min days between changes >= 1   -- now passing
15/15 control(s) remediated and verified.
\end{lstlisting}

\noindent The two \texttt{/etc/login.defs} controls (\texttt{5.5.1.1} and \texttt{5.5.1.2})
\textbf{both} verify, and \texttt{login.defs} appears in the backup tree exactly once ---
confirming the ADR-0006 fix on a real host. After the run, a fresh read-only audit reported
\textbf{100\,\% (32/32)}.

% =========================================================================
\section{Quality, testing \& reproducibility}
\begin{itemize}
  \item \textbf{147 automated tests} (\texttt{pytest}), all OS-independent via \texttt{FakeHost},
        covering each check, the engine, the registry, every reporter, and the remediation
        plan/apply/verify loop --- including the shared-file regression.
  \item \textbf{Deterministic output} ordered by control ID, so report diffs are meaningful.
  \item \textbf{Standard library only} at runtime; \texttt{ruff} for lint/format; Python 3.10+
        to match Ubuntu 22.04's default interpreter.
  \item \textbf{Maintained documentation discipline:} every behaviour change updates
        \texttt{README}, \texttt{CHANGELOG}, the decision log, the control catalogue, and the
        threat model in the same commit.
\end{itemize}

% =========================================================================
\section{Project status}
\textbf{Closed.} The tool audits 32 CIS-mapped controls, remediates the 15 that are safely
auto-fixable, and has been verified before/after on a real Ubuntu 22.04 host (53\,\% $\rightarrow$
100\,\% \emph{of those 32 controls}). The most valuable next step is precisely to grow that
catalogue: the 32 controls are a curated subset of the full benchmark, so broadening coverage ---
and cross-checking results against established scanners such as Lynis and OpenSCAP --- is what
would turn a high score on \emph{this} catalogue into a stronger statement about the host overall.

\vspace{0.6em}
\noindent\textbf{Run it:}
\begin{lstlisting}[style=console]
python -m auditor audit              # read-only audit, console summary
python -m auditor audit --report md  # write a Markdown report
python -m auditor fix --dry-run      # preview changes, touch nothing
python -m auditor fix                # apply safe fixes, with backups
\end{lstlisting}

\vspace{0.4em}
\noindent\small\textcolor{gray}{Generated as the closeout artifact for the Linux Hardening
Auditor portfolio project. Source: \texttt{docs/closeout-report.tex}.}

\end{document}
